\subsection{原函数的存在性}

设 f 是定义在任意区间 \(I \subset R\) 上的函数；为使定义在 I 上的函数 g 是 f 的原函数，必须且只需 g 在每个紧区间 \(J \subset I\) 上的限制是 f 在 J 上的限制的原函数。

定理 1. 设 A 是被滤子 \(\mathfrak{F}\) 滤化的集，\((\mathbf{f}_{\alpha})_{\alpha \in \mathrm{A}}\) 是一族向量函数，取值于 \(\mathbf{R}\) 上的完备赋范空间 E 中，定义在区间 \(\mathbf{I} \subset \mathbf{R}\) 上：对每个 \(\alpha \in \mathrm{A}\)，设 \(\mathbf{g}_{\alpha}\) 是 \(\mathbf{f}_{\alpha}\) 的一个原函数。我们假设：

\(1^{\circ}\) 关于滤子 \(\mathfrak{F}\)，诸函数 \(\mathbf{f}_{\alpha}\) 在 I 的每个紧子集上一致收敛于一个函数 \(\mathbf{f}\)；

\(2^{\circ}\) 存在一点 \(a \in I\)，使得关于滤子 \(\mathfrak{F}\)，族 \((\mathbf{g}_{\alpha}(a))\) 在 E 中有极限。

在这些假设下，诸函数 \(g_{\alpha}\)（关于 F）在 I 的每个紧子集上一致收敛于 f 的一个原函数 g。

根据本小节开头的注记，我们可以只限于 I 是紧区间的情形。

先证明诸 \(g_{\alpha}\) 在 I 上一致收敛于一个连续函数 g。依假设，对每个 \(\varepsilon > 0\) 存在集 \(M \in F\)，使得对属于 M 的任意两个指标 \(\alpha, \beta\)，\(\left\|\mathbf{f}_{\alpha}(x) - \mathbf{f}_{\beta}(x)\right\| \leqslant \varepsilon\) 对每个 \(x \in I\) 成立；因而有（I, p.~15, 定理 2）

\[
\left\| \mathbf {g} _ {\alpha} (x) - \mathbf {g} _ {\beta} (x) - (\mathbf {g} _ {\alpha} (a) - \mathbf {g} _ {\beta} (a)) \right\| \leqslant \varepsilon | x - a | \leqslant \varepsilon l
\]

其中 \(l\) 表示 I 的长度；由于依假设 \(\mathbf{g}_{\alpha}(a)\) 关于 \(\mathfrak{F}\) 趋于一个极限，故由柯西准则，诸 \(\mathbf{g}_{\alpha}\) 在 I 上一致收敛。剩下要证明 \(\mathbf{g}\)（诸 \(\mathbf{g}_{\alpha}\) 的极限）是 \(\mathbf{f}\) 的原函数。

对每个整数 n \textgreater{} 0，设 \(\alpha_{n}\) 是使 \(\left\|\mathbf{f}(x) - \mathbf{f}_{\alpha_{n}}(x)\right\| \leqslant 1/n\) 在 I 上成立的指标；显然序列 \(\left(\mathbf{f}_{\alpha_{n}}\right)\) 一致收敛于 f，并且序列 \(\left(\mathbf{g}_{\alpha_{n}}\right)\) 在 I 上一致收敛于 g。设 \(H_{n}\) 是 I 中使 \(f_{\alpha_{n}}\) 不是 \(g_{\alpha_{n}}\) 的导数的可数子集，并设 H 是诸 \(H_{n}\) 之并，它因而是 I 的可数子集；我们将看到，在每个点 \(x \in \mathrm{I}\)（不属于 \(\mathrm{H}\)）处，函数 \(\mathbf{g}\) 具有等于 \(\mathbf{f}(x)\) 的导数。事实上，与上面一样可以看出，对每个 \(m \geqslant n\) 与每个 \(y \in \mathrm{I}\) 有

\[
\left| \left| \mathbf {g} _ {\alpha_ {m}} (y) - \mathbf {g} _ {\alpha_ {m}} (x) - \left(\mathbf {g} _ {\alpha_ {n}} (y) - \mathbf {g} _ {\alpha_ {n}} (x)\right) \right| \right| \leqslant \frac {2}{n} | y - x |.
\]

令 \(m\) 无限增大，还有

\[
\left| \left| \mathbf {g} (y) - \mathbf {g} (x) - \left(\mathbf {g} _ {\alpha_ {n}} (y) - \mathbf {g} _ {\alpha_ {n}} (x)\right) \right| \right| \leqslant \frac {2}{n} | y - x |
\]

对每个 \(y \in \mathbf{I}\) 成立；而存在 \(h > 0\)，使得对 \(|y - x| \leqslant h\) 与 \(y \in \mathbf{I}\) 有 \(\left\| \mathbf{g}_{\alpha}(y) - \mathbf{g}_{\alpha_n}(x) - \mathbf{f}_{\alpha_n}(x)(y - x) \right\| \leqslant |y - x| / n\)；由于另一方面有 \(\left\| \mathbf{f}(x) - \mathbf{f}_{\alpha_n}(x) \right\| \leqslant 1 / n\)，我们最终得到

\[
\| \mathbf {g} (y) - \mathbf {g} (x) - \mathbf {f} (x) (y - x) \| \leqslant \frac {4}{n} | y - x |
\]

对 \(y \in I\) 与 \(|y - x| \leqslant h\) 成立，这就完成了证明。

推论 1. 在区间 I 上具有原函数的、从 I 到 E 中的映射所成之集 H，是从 I 到 E 中的映射所成的完备向量空间 \(\mathcal{F}_{c}(\mathrm{I};\mathrm{E})\) 的一个闭（从而完备）向量子空间，后者赋予以在 I 的每个紧子集上一致收敛的拓扑（Gen.~Top., X, p.~277）。

推论 2. 设 \(x_0\) 是 I 的一点，对每个函数 \(\mathbf{f} \in \mathcal{H}\)，设 \(\mathrm{P}(\mathbf{f})\) 是 \(\mathbf{f}\) 的在点 \(x_0\) 处为零的原函数；映射 \(\mathbf{f} \mapsto \mathrm{P}(\mathbf{f})\)（从 \(\mathcal{H}\) 到 \(\mathcal{F}_c(\mathrm{I};\mathrm{E})\) 中）是连续线性映射。

定理 1 的推论 1 使我们能按下列程序对某些类函数建立原函数的存在性：若已知属于 \(\mathcal{F}_{c}(\mathrm{I};\mathrm{E})\) 的子集 A 的函数都具有原函数，则属于由 A 生成的向量子空间在 \(\mathcal{F}_{c}(\mathrm{I};\mathrm{E})\) 中的闭包的函数也具有原函数。我们将在下一小节中应用这个方法。

